\chapter*{Обозначения и соглашения}\addcontentsline{toc}{chapter}{Обозначения и соглашения}
Используются СИ и угловые частоты в рад/с, если не сказано иначе. Коэффициенты и неоднозначные буквы определяются локально. Сходство механических, электрических и оптических уравнений не означает равенства размерностей их величин.

\begin{longtable}{p{.21\textwidth}p{.73\textwidth}}
\toprule Символ & Соглашение\\\midrule
$\omega,\nu,f$ & Угловая частота и частота в герцах: $\omega=2\pi\nu=2\pi f$.\\
$\gamma,Q$ & В $\ddot x+2\gamma\dot x+\omega_0^2x=0$ амплитуда убывает с декрементом $\gamma$, энергия — $2\gamma$; $Q\simeq\omega_0/(2\gamma)$.\\
$\tau_p,\kappa$ & Время жизни энергии/фотонов и скорость их потерь: $\kappa=1/\tau_p$, $Q=\omega\tau_p$. Амплитудный декремент $\kappa/2$.\\
$A,\phi,\Delta$ & Амплитуда, фаза и расстройка. Определение вращающейся системы и знак расстройки указываются локально.\\
$\langle\cdot\rangle$ & В колебательных задачах — среднее по быстрому периоду при фиксированных медленных переменных; в статистических — ансамблевое среднее.\\
\bottomrule
\end{longtable}

\paragraph{Преобразование Фурье и спектры.}
\begin{equation*}
 X(\omega)=\int_{-\infty}^{\infty}x(t)\e^{\ii\omega t}\dd t,\qquad
 x(t)=\frac1{2\pi}\int_{-\infty}^{\infty}X(\omega)\e^{-\ii\omega t}\dd\omega.
\end{equation*}
Плоская волна записана как $\operatorname{Re}[A\e^{\ii\bm k\cdot\bm r-\ii\omega t}]$; у пассивной волны вдоль положительной оси $\operatorname{Im}k>0$. При другом явно объявленном соглашении знаки комплексных коэффициентов меняются совместно.

Для центрированного вещественного стационарного процесса $C(\tau)=\langle x(t+\tau)x(t)\rangle$, $S(\omega)=\int C(\tau)\e^{\ii\omega\tau}\dd\tau$ — двусторонняя плотность:
\begin{equation*}
 \langle x^2\rangle=\frac1{2\pi}\int_{-\infty}^{\infty}S(\omega)\dd\omega,
 \quad G(f)=2S(2\pi f)\ (f>0),\quad
 \langle x^2\rangle=\int_0^\infty G(f)\dd f.
\end{equation*}
$G$ — односторонняя плотность в герцах. Среднее процесса вычитают; иначе при нулевой частоте появляется дельта-линия. Квантовые симметризованные и несимметризованные спектры различаются отдельно.

\paragraph{Амплитуда, мощность и ширина.}
Комплексные гармонические амплитуды — пиковые, если не сказано иное; средняя мощность содержит множитель $1/2$. Поле, интенсивность и число фотонов не взаимозаменяемы. При $W\propto\e^{-t/\tau_p}$ полная ширина пассивного резонанса мощности в герцах равна $1/(2\pi\tau_p)$; лазерная шумовая линия определяется фазовой диффузией. Полную ширину и полуширину, интенсивностные и амплитудные потери сравнивают только после перевода нормировок.
